\documentclass[a4paper]{article}

% Basic packages
\usepackage[fontset=fandol]{ctex}
\usepackage{amsmath, amssymb}
\usepackage{graphicx}
\usepackage[margin=2.5cm]{geometry}
\usepackage[most]{tcolorbox}
\usepackage{etoolbox}
\usepackage{listings}
\usepackage{hyperref}
\usepackage{booktabs}
\usepackage{subcaption}
\usepackage{float}
\usepackage{tikz}
\IfFileExists{pgfplots.sty}{
    \usepackage{pgfplots}
    \pgfplotsset{compat=1.18}
}{}

% Highlight box palette: low-saturation and print-friendly.
% Visual weight is ordered by teaching role, so a page mixing several boxes still
% reads important > warning > knowledge > dialogue instead of competing for attention.
% Title text is white on every title bar; contrast stays within WCAG AA (4.5:1--6.0:1).
\definecolor{kbBack}{HTML}{F2F6FA}\definecolor{kbFrame}{HTML}{9DB4CC}\definecolor{kbTitle}{HTML}{52708F}
\definecolor{ibBack}{HTML}{FBF5E7}\definecolor{ibFrame}{HTML}{C9A961}\definecolor{ibTitle}{HTML}{7D5F1E}
\definecolor{wbBack}{HTML}{FCF2F0}\definecolor{wbFrame}{HTML}{D6A096}\definecolor{wbTitle}{HTML}{9E4F43}
\definecolor{dbBack}{HTML}{F4F7F3}\definecolor{dbFrame}{HTML}{AEC2AC}\definecolor{dbTitle}{HTML}{5F7F64}

\tcbset{softbox/.style={
    enhanced, boxrule=0.8pt, arc=2pt, left=3mm, right=3mm, top=2mm, bottom=2mm,
    coltitle=white, fonttitle=\bfseries, toptitle=0.6mm, bottomtitle=0.6mm
}}

% Highlight boxes
\newtcolorbox{knowledgebox}[1]{
    softbox, colback=kbBack, colframe=kbFrame, colbacktitle=kbTitle, title=#1
}

\newtcolorbox{importantbox}[1]{
    softbox, colback=ibBack, colframe=ibFrame, colbacktitle=ibTitle, title=#1
}

\newtcolorbox{warningbox}[1]{
    softbox, colback=wbBack, colframe=wbFrame, colbacktitle=wbTitle, title=#1
}

\newtcolorbox{dialoguebox}[1]{
    softbox, breakable, colback=dbBack, colframe=dbFrame, colbacktitle=dbTitle, title=#1
}

% Listings
\lstset{
    language=Python,
    basicstyle=\ttfamily\small,
    keywordstyle=\color{blue},
    stringstyle=\color{red!60!black},
    commentstyle=\color{green!60!black},
    breaklines=true,
    frame=single,
    numbers=left,
    numberstyle=\tiny\color{gray},
    captionpos=b,
    extendedchars=false
}

% Front-page metadata
\newcommand{\notetitle}{课程笔记：[在此填写课程主题]}
\newcommand{\noteauthors}{五道口纳什 \& Codex}
\newcommand{\notedate}{\today}
\newcommand{\videochannel}{[在此填写频道或作者]}
\newcommand{\videopublishdate}{[在此填写发布日期]}
\newcommand{\videoduration}{[在此填写视频时长]}
\newcommand{\videourl}{}
\newcommand{\videocoverpath}{}

\begin{document}

\begin{titlepage}
\centering
\vspace{1.2cm}
{\huge\bfseries \notetitle\par}
\vspace{0.8cm}
{\large \noteauthors\par}
\vspace{0.3cm}
{\large \notedate\par}
\vspace{1.2cm}

\ifdefempty{\videocoverpath}{
{\small 请在生成笔记时填入视频封面图的本地路径。\par}
}{
\includegraphics[width=0.82\textwidth,height=0.45\textheight,keepaspectratio]{\videocoverpath}\par
}

\vfill
\begin{tcolorbox}[width=0.9\textwidth, colback=black!2!white, colframe=black!60, sharp corners]
\textbf{视频作者/频道}：\videochannel\par
\textbf{发布日期}：\videopublishdate\par
\textbf{视频时长}：\videoduration\par
\textbf{视频链接}：
\ifdefempty{\videourl}{
未填写
}{
\href{\videourl}{\nolinkurl{\videourl}}
}
\par\textbf{PDF制作Skill}：\href{https://github.com/wdkns/wdkns-skills}{\nolinkurl{https://github.com/wdkns/wdkns-skills}}
\end{tcolorbox}
\end{titlepage}

\tableofcontents
\newpage

%% --- 正文内容开始 --- %%

% Replace this block with the generated notes.
% Fill the metadata block above, especially \videocoverpath, so the front page
% contains the original video cover image.
% Keep images outside knowledgebox/importantbox/warningbox/dialoguebox.
% For interview, panel, podcast, or conversation videos, optionally use dialoguebox
% to preserve a short high-signal original dialogue segment that is especially
% informative, vivid, funny, or intuitive. It may be a single exchange or several
% tightly connected turns. Keep speaker labels and source timestamps; do not dump
% long transcript blocks.
% \begin{dialoguebox}{原始对话片段（00:18:42--00:19:05）}
% \textbf{主持人}：这里放一句短问题或追问。
%
% \textbf{嘉宾}：这里放一句真正有解释力、临场感或直觉价值的回答。
%
% \textbf{主持人}：如有必要，可以保留紧接着的多轮澄清、反驳或补充。
% \end{dialoguebox}
% For any figure that comes from a video frame or a crop of a frame,
% include the source time interval as a bottom footnote on the same page.
% A stable pattern is:
% \begin{figure}[H]
% \centering
% \includegraphics[width=\textwidth]{...}
% \caption{... \protect\footnotemark}
% \end{figure}
% \footnotetext{视频画面时间区间：00:12:31--00:12:46。若为裁剪图，时间区间仍对应原视频画面。}
% End the document with a final section such as:
% \section{总结与延伸}
% covering the speaker's closing discussion and your own synthesis.

%% --- 正文内容结束 --- %%

\end{document}
